\documentclass{article}

\usepackage{booktabs}
\usepackage{tabularx}

\title{Development Plan\\\progname}

\author{\authname}

\date{}

\input{../Comments.text}
\input{../Common.text}

\begin{document}

\maketitle

\begin{table}[hp]
\caption{Revision History} \label{TblRevisionHistory}
\begin{tabularx}{\textwidth}{llX}
\toprule
\textbf{Date} & \textbf{Developer(s)} & \textbf{Change}\\
\midrule
September 23 & Tyler Yue & Added rough draft of workflow plan, communication plan and team meeting plan.\\
September 25 & Qusay Qadir & Use of AI, In Case of Emergency, Expected Technology. \\
... & ... & ...\\
\bottomrule
\end{tabularx}
\end{table}

\newpage{}

\wss{Put your introductory blurb here.  Often the blurb is a brief roadmap of
what is contained in the report.}

\wss{Additional information on the development plan can be found in the
\href{https://gitlab.cas.mcmaster.ca/courses/capstone/-/blob/main/Lectures/L02b_POCAndDevPlan/POCAndDevPlan.pdf?ref_type=heads}
{lecture slides}.}

\section{Confidential Information?}

\wss{State whether your project has confidential information from industry, or
not.  If there is confidential information, point to the agreement you have in
place.}

\wss{For most teams this section will just state that there is no confidential
information to protect.}
\section{IP to Protect}

\wss{State whether there is IP to protect.  If there is, point to the agreement.
All students who are working on a project that requires an IP agreement are also
required to sign the ``Intellectual Property Guide Acknowledgement.''}

\section{Copyright License}

\wss{What copyright license is your team adopting.  Point to the license in your
repo.}

\section{Team Meeting Plan}

\wss{This section is for meetings between the team members, and meetings between 
the team and their supervisor/advisor (as appropriate).}

\wss{How often will the team meet? where? when?}

\wss{If the meeting is a physical location (not virtual), out of an abundance of
caution for safety reasons you shouldn't put the location online}

\wss{To help schedule meetings, you can use the lecture/tutorial time when we do 
not have anything else scheduled.}

\wss{How often will you meet with your supervisor/advisor?  when?  where?}

\wss{Will meetings be virtual?  At least some meetings should likely be
in-person.}

\wss{How will the meetings be structured?  There should be a chair for all meetings.  
There should be an agenda for all meetings.}

\section{Team Communication Plan}

\wss{Issues on GitHub should be part of your communication plan.}

\section{Team Member Roles}

\wss{You should identify the types of roles you anticipate, like notetaker,
leader, meeting chair, reviewer.  Assigning specific people to those roles is
not necessary at this stage.  In a student team the role of the individuals will
likely change throughout the year.}

\section{Workflow Plan}

\begin{itemize}
	\item How will you be using git, including branches, pull request, etc.?
	\item How will you be managing issues, including template issues, issue
	classification, etc.?
  \item Use of CI/CD
\end{itemize}

\section{Project Decomposition and Scheduling}

\begin{itemize}
  \item How will you be using GitHub projects?
  \item Include a link to your GitHub project
\end{itemize}

\wss{How will the project be scheduled?  This is the big picture schedule, not
details. You will need to reproduce information that is in the course outline
for deadlines.}

\section{Proof of Concept Demonstration Plan}

The primary risk to the success of our project is whether we are able to construct and maintain a model that accurately represents the relationships between a software system’s implementation and engineering intent.

Our tool is dependent on understanding how files, classes, code, and other artifacts relate or map to requirements, constraints, and design decisions. We must be able to justify each of these relationships and update each of these relationships as code is added or updated, and new features are added. This functionality is required before we can focus on higher level, user facing functionality, such as constraint aware code generation, design rational tracking, etc.

The main challenges in demonstrating this representation will be determining how these relationships can be identified accurately, how these relationships can be justified, and how the model can remain synchronized as the software project changes. A requirement, design decision, or new feature can affect multiple files, classes, and code, making these relationships difficult to identify and maintain. As the implementation evolves, previously established relationships may become outdated or invalid, requiring the model to detect and reflect these changes while preserving accurate traceability.

Therefore, for our proof of concept, we will focus on demonstrating this underlying representation, rather than showcasing a subset of features that we aim to eventually have, as this piece of functionality is the foundation for any features we will add. Given a software project and a set of requirements, constraints, and design choices, the system will construct a model that maps these to implementations. For each of these relationships, the system should be able to provide a justification of how the connection was established.


\section{Expected Technology}

% \wss{What programming language or languages do you expect to use?  What external
% libraries?  What frameworks?  What technologies.  Are there major components of
% the implementation that you expect you will implement, despite the existence of
% libraries that provide the required functionality.  For projects with machine
% learning, will you use pre-trained models, or be training your own model?  }

% \wss{The implementation decisions can, and likely will, change over the course
% of the project.  The initial documentation should be written in an abstract way;
% it should be agnostic of the implementation choices, unless the implementation
% choices are project constraints.  However, recording our initial thoughts on
% implementation helps understand the challenge level and feasibility of a
% project.  It may also help with early identification of areas where project
% members will need to augment their training.}

\begin{table}[h]
  \centering
  \small
  \renewcommand{\arraystretch}{1.3}
  \begin{tabular}{|p{0.15\linewidth}|p{0.2\linewidth}|p{0.1\linewidth}|p{0.4\linewidth}|}
  \hline
  \textbf{Category} & \textbf{Technology} & \textbf{Version} & \textbf{Purpose} \\
  \hline
  Language & Python & 3.14 & Backend services, intent model processing, and code analysis \\
  \hline
  Language & TypeScript & 7.0 & VS Code extension and user interface \\
  \hline
  Runtime & Node.js & 24 LTS & Build tooling and extension development \\
  \hline
  Framework & VS Code Extension API & 1.139 & Integration with the editor, commands, panels, and file watching \\
  \hline
  Framework & React & 19.3 & Webview user interface and diagram rendering \\
  \hline
  Library & MCP Python SDK & 2.2 & Exposing the intent model to AI coding assistants \\
  \hline
  Framework & LangGraph & 1.2 & Agent orchestration with human-in-the-loop approval \\
  \hline
  Tool & LangSmith & 0.14 & Tracing and evaluating agent and LLM behaviour \\
  \hline
  Version control & Git & 2.55 & Source control and versioning of the intent model \\
  \hline
  Platform & GitHub and GitHub Actions & N/A & Repository hosting, continuous integration, and pull request checks \\
  \hline

\end{tabular}
\caption{Initial implementation technologies}
\label{tab:tech-stack}
\end{table}

\paragraph{Components implemented by the team.} Although libraries exist for parsing, 
graph rendering, and LLM orchestration, the team expects to implement the core of the 
system itself: the intent modelling language and its schema, the pipeline that extracts 
proposed model changes from natural-language conversation, the traceability graph 
linking intent to UML models, code, and tests, and the logic that compiles constraints 
into executable checks. No existing tool provides these capabilities in the form the 
project requires.


\paragraph{Topics left to discuss:}

\begin{itemize}
\item Specific libraries
\item LLM Pre-trained models
\item Specific linter tool (if appropriate)
\item Specific unit testing framework
\item Investigation of code coverage measuring tools
\item Specific performance measuring tools (like Valgrind), if
\end{itemize}

% \wss{git, GitHub and GitHub projects should be part of your technology.}

\section{Use of AI and LLMs}

% \wss{This section is a continuation of the previous section.  The role of LLMs
% in your project as important enough to warrant their own section.  
% How will your team be using LLMs for code and documentation development?  
% What tools will you be using?  How will you verify the output of your LLMs?}

For the project the model artifact that we hope to generate is based on the input provided by the user in natural-language to an AI Assistant. The AI would also be used to generate code based on the intent model, test cases, class and code structure descriptions, explanations, and relationships with the rest of the code, and comments on PRs ( if we get to one of our stretch goals). 
We would use AI to help assist in code generation after we design and engineer our intent model. Lastly, LLMs would be used to help proof read, edit, and mock evaluate our documentation based on the rubric to catch details that might have been missed. 
For verification there would need to be a human-in-the-loop, this means that anything that the LLM generates should be used or pushed to our codebase until a member of the team can explicitly explain its use, function, and implementation. 

\section{Coding Standard}

The project will adopt a consistent coding standard focused on readability, maintainability, and collaboration. Each language used will follow its established conventions for naming, formatting, type usage, and code organization, with automated linting and formatting tools used where appropriate.

Functionality developed as part of the project will include accompanying documentation describing its purpose, behaviour, and usage. Documentation will be organized consistently and updated alongside the code as the project evolves.


\newpage{}

\section{Appendix --- Generative AI Use Disclosure}

\wss{ChatGPT (version 5) was used to brainstorm ideas for the template for this disclose section.}

\wss{This section is about the use of AI to write this document, not your
planned use of AI for your project.  You discuss that topic in one of the
sections in the body of this report.}

\subsection{AI Tools Used}

\wss{Give the name of the AI tool(s) used and the version.}

\subsection{How and Where Used}

\wss{Explain what the tool(s) were used for.  Examples include: brainstorming,
explaining a technical concept, reviewing the document, generating or modifying
text, generating or modifying diagrams, identifying errors or omissions or
inconsistencies.}

\wss{For each of the uses, identify which parts of the document were affected.}

\wss{You do not need to report every interaction, but you should disclose the 
uses that materially contributed to your work}

\subsection{Verification of AI-Generated Content}

\wss{\begin{itemize}
    \item Describe how the team checked AI-generated or AI-assisted material
    before including it in the document.
    \item In particular, verify factual claims, technical assertions,
    requirements, calculations, references, and other information for which an
    AI system may produce incorrect or fabricated results.
    \item \textbf{The team is responsible for the accuracy, completeness,
consistency, and appropriateness of the submitted document, regardless of
whether material was generated by a human or an AI system.}
\end{itemize}}

\newpage{}

\section*{Appendix --- Reflection}

\wss{Not required for CAS 741}

\input{../Reflection.text}

\begin{enumerate}
    \item Why is it important to create a development plan prior to starting the
    project?
    \item In your opinion, what are the advantages and disadvantages of using
    CI/CD?
    \item What disagreements did your group have in this deliverable, if any,
    and how did you resolve them?
\end{enumerate}

\newpage{}

\section*{Appendix --- Team Charter}

\wss{borrows from
\href{https://engineering.up.edu/industry_partnerships/files/team-charter.pdf}
{University of Portland Team Charter}}

\subsection*{External Goals}

\wss{What are your team's external goals for this project? These are not the
goals related to the functionality or quality fo the project.  These are the
goals on what the team wishes to achieve with the project.  Potential goals are
to win a prize at the Capstone EXPO, or to have something to talk about in
interviews, or to get an A+, etc.}

\subsection*{Attendance}

\subsubsection*{Expectations}

\wss{What are your team's expectations regarding meeting attendance (being on
time, leaving early, missing meetings, etc.)?}

\subsubsection*{Acceptable Excuse}

\wss{What constitutes an acceptable excuse for missing a meeting or a deadline?
What types of excuses will not be considered acceptable?}

\subsubsection*{In Case of Emergency}

% \wss{What process will team members follow if they have an emergency and cannot
% attend a team meeting or complete their individual work promised for a team
% deliverable?}

If a team member has an emergency, they will notify the team as soon as reasonably possible through the team's main communication channel. 
Where possible, they will share the status of their current work, including what is complete, what remains, and where the work is located.

A delegate will then be assigned to take over the absent member's action items for that week. If the absent member is unable to 
arrange this themselves, the remaining team members will assign a delegate at the next meeting or through the team channel. 
If the delegate is already overloaded, the team will hold a meeting to create a remediation plan and redistribute the work among available members. 

To make transfers of responsibility seamless, all members should know what everyone else is working on through the weekly meetings, 
even if not the implementation details. Because the team follows shared coding standards and keeps work in the shared repository, 
any member should be able to understand, navigate, and pick up another member's work based on prior knowledge.

\subsection*{Accountability and Teamwork}

\subsubsection*{Quality} 

\wss{What are your team's expectations regarding the quality
of team members' preparation for team meetings and the quality of the
deliverables that members bring to the team?}

\subsubsection*{Attitude}

\wss{What are your team's expectations regarding team members' ideas,
interactions with the team, cooperation, attitudes, and anything else regarding
team member contributions?  Do you want to introduce a code of conduct?  Do you
want a conflict resolution plan?  Can adopt existing codes of conduct.}

\subsubsection*{Stay on Track}

\wss{What methods will be used to keep the team on track? How will your team
ensure that members contribute as expected to the team and that the team
performs as expected? How will your team reward members who do well and manage
members whose performance is below expectations?  What are the consequences for
someone not contributing their fair share?}

\wss{You may wish to use the project management metrics collected for the TA and
instructor for this.}

\wss{You can set target metrics for attendance, commits, etc.  What are the
consequences if someone doesn't hit their targets?  Do they need to bring the
coffee to the next team meeting?  Does the team need to make an appointment with
their TA, or the instructor?  Are there incentives for reaching targets early?}

\subsubsection*{Team Building}

Team cohesion will be built through regular meetings, frequent communication, and collaborative development practices. Each merge request will be reviewed by at least two other team members, ensuring that everyone remains aware of the work being completed across the project, frequent contact between team members, and that team members are not working in isolation.

Higher-level design and project decisions will also be discussed as a group during team meetings.

Outside of development work, the team may also organize informal activities, such as meeting for coffee or spending time together outside scheduled project meetings, to strengthen communication and working relationships which will help enhance team cohesion.

\subsubsection*{Decision Making} 

Smaller implementation level decisions will follow the coding standards, naming conventions, and design patterns established by the team beforehand. Within these guidelines, team members will have discretion over how they implement their assigned work. Any concerns with an individual member's implementation level decisions can be raised during merge request reviews and can be raised to a team discussion which can also set future precedence.

Larger decisions that affect the overall project direction, architecture, or design will be discussed as a team during meetings prior to starting any of the implementation. The team will aim to reach a consensus, however, if not possible, the minimum expectation is that all members can understand and agree with the reasoning behind the final decision, even if their preferences differ.

\end{document}